% article-template.tex — Template for a research paper (amsart class)
%
% Compile: latexmk (uses .latexmkrc in this directory)
% Engine:  XeLaTeX

\documentclass[12pt]{amsart}

% ─── Encoding and Fonts ─────────────────────────────────────────────────────
\usepackage{fontspec}
\usepackage{microtype}

% ─── Mathematics (loaded by amsart, but we add extras) ───────────────────────
\usepackage{mathtools}
\usepackage{amssymb}

% ─── Page Layout ─────────────────────────────────────────────────────────────
\usepackage{geometry}
\geometry{margin=1in}

% ─── Theorem-like Environments ───────────────────────────────────────────────
% amsart provides \newtheorem; we define our environments here.
\theoremstyle{plain}
\newtheorem{theorem}{Theorem}[section]
\newtheorem{lemma}[theorem]{Lemma}
\newtheorem{proposition}[theorem]{Proposition}
\newtheorem{corollary}[theorem]{Corollary}
\newtheorem{conjecture}[theorem]{Conjecture}

\theoremstyle{definition}
\newtheorem{definition}[theorem]{Definition}
\newtheorem{example}[theorem]{Example}

\theoremstyle{remark}
\newtheorem{remark}[theorem]{Remark}

% ─── Graphics, Diagrams, and Colors ─────────────────────────────────────────
\usepackage{graphicx}
\usepackage{xcolor}
\usepackage{tikz}
\usepackage{tikz-cd}

% ─── Cross-referencing ───────────────────────────────────────────────────────
\usepackage[
  colorlinks=true,
  linkcolor=blue!40!black,
  citecolor=green!40!black,
  urlcolor=magenta!40!black
]{hyperref}
\usepackage[capitalise,noabbrev]{cleveref}

% ─── Bibliography ────────────────────────────────────────────────────────────
\usepackage[
  backend=biber,
  style=alphabetic,
  sorting=nyt,
  maxbibnames=10
]{biblatex}
\addbibresource{references.bib}

% ─── Collaboration Tools ────────────────────────────────────────────────────
\usepackage[colorinlistoftodos,textsize=small]{todonotes}

% ─── Custom Macros ───────────────────────────────────────────────────────────

% Number sets
\newcommand{\NN}{\mathbb{N}}
\newcommand{\ZZ}{\mathbb{Z}}
\newcommand{\QQ}{\mathbb{Q}}
\newcommand{\RR}{\mathbb{R}}
\newcommand{\CC}{\mathbb{C}}

% Delimiters
\DeclarePairedDelimiter{\abs}{\lvert}{\rvert}
\DeclarePairedDelimiter{\norm}{\lVert}{\rVert}
\DeclarePairedDelimiter{\inner}{\langle}{\rangle}

% Operators
\DeclareMathOperator{\id}{id}

% ─── Document Metadata ──────────────────────────────────────────────────────
\title{Title of the Paper}

\author{Author One}
\address{Department of Mathematics, University of Example, City, State, ZIP}
\email{author.one@example.edu}

\author{Author Two}
\address{Department of Mathematics, Another University, City, State, ZIP}
\email{author.two@example.edu}

\date{\today}

\subjclass[2020]{Primary XXXXX; Secondary YYYYY}
\keywords{keyword one, keyword two, keyword three}

\begin{abstract}
  This paper investigates the properties of widgets
  and establishes a characterization of compact widgets
  in terms of sequential compactness.
  We prove that every compact widget is sequentially compact
  and discuss the converse under additional hypotheses.
\end{abstract}

% ═══════════════════════════════════════════════════════════════════════════════
\begin{document}

\maketitle

% ─── Introduction ────────────────────────────────────────────────────────────
\section{Introduction}\label{sec:introduction}

The study of widgets was initiated by Doe and Smith \autocite{placeholder2025}.
In this paper, we establish the following result.

\begin{theorem}\label{thm:main}
  Every compact widget is sequentially compact.
\end{theorem}

The proof of \cref{thm:main} is given in \cref{sec:proof}.

\subsection*{Organization}

In \cref{sec:preliminaries}, we recall the necessary definitions.
\Cref{sec:proof} contains the proof of \cref{thm:main}.
We conclude in \cref{sec:remarks} with open questions.

\subsection*{Acknowledgements}

\todo{Add acknowledgements.}

% ─── Preliminaries ──────────────────────────────────────────────────────────
\section{Preliminaries}\label{sec:preliminaries}

We fix notation and recall the basic definitions.

\begin{definition}\label{def:widget}
  A \emph{widget} is a pair $(X, \tau)$ where $X$ is a set and
  $\tau \subseteq \mathcal{P}(X)$ satisfies:
  \begin{enumerate}
    \item $\emptyset, X \in \tau$;
    \item $\tau$ is closed under arbitrary unions;
    \item $\tau$ is closed under finite intersections.
  \end{enumerate}
\end{definition}

\begin{definition}\label{def:compact}
  A widget $(X, \tau)$ is \emph{compact} if every open cover of $X$ has a
  finite subcover.
\end{definition}

% ─── Main Proof ──────────────────────────────────────────────────────────────
\section{Proof of the Main Theorem}\label{sec:proof}

The key ingredient is the following lemma.

\begin{lemma}\label{lem:subsequence}
  Let $(X, \tau)$ be a compact widget and let $(x_n)$ be a sequence in $X$.
  Then $(x_n)$ has a convergent subsequence.
\end{lemma}

\begin{proof}
  Suppose for contradiction that no subsequence converges.
  For each $x \in X$, choose an open neighborhood $U_x$ containing
  only finitely many terms of the sequence.
  Then $\{U_x\}_{x \in X}$ is an open cover with no finite subcover,
  contradicting compactness.
\end{proof}

\begin{proof}[Proof of \cref{thm:main}]
  This follows immediately from \cref{lem:subsequence}.
\end{proof}

% ─── Concluding Remarks ─────────────────────────────────────────────────────
\section{Concluding Remarks}\label{sec:remarks}

\begin{conjecture}\label{conj:converse}
  Every sequentially compact widget is compact.
\end{conjecture}

\Cref{conj:converse} is known to hold under additional set-theoretic
hypotheses; see \autocite{example2024}.

% ─── Bibliography ────────────────────────────────────────────────────────────
\printbibliography

\end{document}

%%% Local Variables:
%%% mode: LaTeX
%%% TeX-master: t
%%% End:
